\lecture{3}
\title{Approximation Theory}

\begin{document}

These are the important questions we will be answering
during this lecture.
\begin{itemize}
    \item \textbf{Approximation.}
    What is the best $\hat{\mathcal{R}}(\theta^*)$
    (empirical risk) we can achieve with our model?
    \item \textbf{Optimization.}
    How good is our learning algorithm
    at minimizing $\hat{\mathcal{R}}(\theta)$ (empirical risk)?
    \item \textbf{Generalization.}
    How different is $\hat{\mathcal{R}}(\theta)$ (empirical risk)
    from $\mathcal{R}(\theta)$ (true risk)?
\end{itemize}

\dots and I realize that I have no reason to be here,
because this lecture covers the exact same content as:
\begin{itemize}
    \item The ``Neural Network Universal Approximation'' section
    of \href{/6.790/lectures/12/#neural-network-universal-approximation}{Lecture 12}
    of 6.790.
    \item The ``Neural Network Depth'' section
    of \href{/6.790/lectures/13/#neural-network-depth}{Lecture 13} of 6.790.
    \item The ``Estimation Error Bounding'' section
    of \href{/6.790/lectures/10/#estimation-error-bounding}{Lecture 10}
    of 6.790.
\end{itemize}

This class offers a little more detail on the
``Universal Approximation'' theorem;
it provides a (very non-rigorous) proof sketch of:

\begin{theorem}{Universal Approximation}
    Suppose $g \colon [0, 1]^d \to \mathbb{R}$ is an $L$-Lipschitz function.
    Then for any $\epsilon > 0$,
    there exists a 3-layer ReLU network $f$
    with $O((L/\epsilon)^d)$ units such that:
    \[ \int_{[0, 1]^d} |f(\mathbf{x}) - g(\mathbf{x})|
    \, \mathrm{d}\mathbf{x} \leq 2 \epsilon. \]
\end{theorem}

Actually, no, turns out 6.790 said this already, too.
I'm leaving.

\end{document}